\section{A Banach space of differences of strong energy potentials}
\label{sec:r12-energy-difference}

\paragraph{Exact standing hypotheses and frozen inputs.}
Let $K/\mathbf Q_p$ be the original finite extension with its original
absolute value. Let $X$ be good, Hausdorff, boundaryless, quasi-smooth,
and pure of dimension $n$. Choose a compact original affinoid domain
$B\subset X$, functions $f_a,g_b\in\mathcal O(B)$ and positive real
thresholds $R_a,r_b$, such that
\[
 K_0=\{|f_a|\leq R_a,\ |g_b|\geq r_b\}
                         \Subset\operatorname{Int}_X(B),\qquad
 V=\{|f_a|<R_a,\ |g_b|>r_b\}\ne\varnothing .
 \tag{ED1}
\]
Let $S$ be an arbitrary original symmetric CLD middle-positive,
separately $d'$- and $d''$-closed $(p,p)$-current on an ordinary
neighborhood of $K_0$. Put $q=n-p\geq1$ and $d=q+1$.
No strong positivity, geometric representation or approximation of
$S$ is assumed. Every current below is tested on the full original
CLD test sheaf, with every actual moment refinement.

Write $\mathscr C=\widehat{\mathcal E}^{1}_{\rm str}(V)$ for the
precise R11 strong completion, with its encoded approximation
boundary condition. Its elements are nonpositive Borel potentials
modulo local ordinary-open outer $S$-capacity zero. We use the R11
inputs (GE10)--(GE18), (IC10)--(IC22), (CC1)--(CC11), (CI1)--(CI8)
and the completed extension of (EV1). They provide, respectively,
the symmetric global mixed energy, its positive energy matrices,
the complete strong cone, its intrinsic original currents, and its
monotone closure. All these statements are \texttt{INTERNAL\_INPUT}
in their printed R11 scopes.

For $u_0,\ldots,u_q\in\mathscr C$, abbreviate
\[
 B(u_0,\ldots,u_q)=\int_V(-u_0)\,\widehat N_q(u_1,\ldots,u_q),
 \quad e(u)=B(u^{[d]}),\quad r(u)=e(u)^{1/d}.
 \tag{ED2}
\]
The frozen inequalities are $B(\mathbf u)\leq\prod_i r(u_i)$ and
$r(u+v)\leq r(u)+r(v)$. The zero-energy class is zero. Completed
potentials are finite outside outer-null sets; all differences below
are formed there and assigned any finite Borel value on the joint
exceptional set. In particular no expression $(-\infty)-(-\infty)$
is used as a pointwise number.

\begin{definition}[Energy-difference space and its norm]
\label{def:r12-energy-difference}
Let $\mathbb D_S=\mathscr C-\mathscr C$ be the real vector space of
such Borel differences modulo local outer-capacity equality. Define
\[
 \|h\|_{\rm ed}
   =\inf\{r(a)+r(b):a,b\in\mathscr C,\ h=a-b\}.
 \tag{ED3}
\]
This is a norm defined by the energy cone, not the discrepancy $I$.
It is not declared to be the ordinary supremum norm or a Hilbert norm.
\end{definition}

\begin{lemma}[Continuous extension of the mixed energy]
\label{lem:r12-difference-mixed-extension}
There is a unique symmetric real $d$-linear form on $\mathbb D_S$
extending $B$, and it satisfies
\[
       |B(h_0,\ldots,h_q)|\leq\prod_{i=0}^q\|h_i\|_{\rm ed}.
 \tag{ED4}
\]
Moreover $\|u\|_{\rm ed}=r(u)$ for every $u\in\mathscr C$.
\end{lemma}

\noindent\textbf{Status: \texttt{NEW\_PROVED}} for the full
unbounded cone-difference construction; the mixed energy bound on
the positive cone is the explicit frozen input.

\begin{proof}
Expand each difference into its two cone terms and extend by the
finite multilinear expansion. If $a-b=a'-b'$, then
$a+b'=a'+b$ in the original relative cone. Separate additivity of
the frozen $B$ therefore makes the expansion independent of that
presentation, one slot at a time. Positive real homogeneity and
changing the signs of differences give real multilinearity.
For specified decompositions $h_i=a_i-b_i$, the triangle inequality
and cone Holder give the bound $\prod_i(r(a_i)+r(b_i))$.
Taking independent infima proves (ED4).

The definition gives $\|u\|_{\rm ed}\leq r(u)$. Applying (ED4)
to the diagonal tuple gives $e(u)\leq\|u\|_{\rm ed}^{d}$, so the
reverse inequality follows. This argument also covers zero energy.
\end{proof}

\begin{theorem}[Actual Borel realization and Banach completeness]
\label{thm:r12-energy-difference-Banach}
Formula (ED3) is a norm, $\mathbb D_S$ is complete for that norm,
and $\mathscr C$ is a closed convex cone in this Banach space.
Norm convergence implies convergence in the full local outer
$S$-capacity topology. More precisely, with
$C_0=d2^d$, one has
\[
 C_S^*(\{|h|>t\};V)\leq C_0 t^{-d}\|h\|_{\rm ed}^{d}
                    \qquad(h\in\mathbb D_S,\ t>0).
 \tag{ED5}
\]
For $u,v\in\mathscr C$,
\[
 I(u,v)\leq\bigl(r(u)^q+r(v)^q\bigr)\|u-v\|_{\rm ed}.
 \tag{ED6}
\]
The result does not identify this norm topology with the strong
discrepancy topology on $\mathscr C$.
\end{theorem}

\noindent\textbf{Status: \texttt{NEW\_PROVED}}.

\begin{proof}
The infimum in (ED3) is finite. Minkowski on $\mathscr C$ proves
the triangle inequality; switching the two cone terms proves
symmetry, and homogeneity follows by rescaling decompositions.
The R11 capacity integral bound $J_1^*(u)\leq2^d e(u)$ and
monotonicity of the sublevels give
\[
 \frac{t^d}{d}C_S^*(\{u<-t\};V)
       \leq\int_0^t s^{d-1}C_S^*(\{u<-s\};V)\,ds
       \leq2^d e(u).
\]
For $h=a-b$, the set $\{|h|>t\}$ is contained, modulo an
outer-null set, in $\{a<-t\}\cup\{b<-t\}$. Hence its outer
capacity is at most $C_0t^{-d}(r(a)^d+r(b)^d)$, bounded by
$C_0t^{-d}(r(a)+r(b))^d$. Infimizing proves (ED5), and shows
that a zero seminorm has the zero relative Borel class. Local
outer-null sets here also have zero outer capacity in the fixed
domain: use a countable compact exhaustion and the R11 domain
comparison on its cores. The same comparison transfers norm
convergence to every smaller-domain original capacity.

We prove completeness inside the space of actual Borel differences,
rather than introducing an abstract completion with unidentified
elements. For a norm-Cauchy sequence choose a subsequence $(h_j)$
whose successive increments have decompositions
\[
 h_{j+1}-h_j=a_j-b_j,\qquad
 \sum_{j\geq1}(r(a_j)+r(b_j))<\infty.
 \tag{ED7}
\]
Write $h_1=a_0-b_0$. The partial sums
$A_N=a_0+\sum_{j=1}^N a_j$ and
$B_N=b_0+\sum_{j=1}^N b_j$ decrease as Borel potentials,
outside one countable union of outer-null sets. Their energy
roots are bounded by the corresponding sums of roots, so the
R11 decreasing-closure theorem (CI8) gives strong limits
$A,B\in\mathscr C$. Apply the same theorem to each tail series.
It yields $A-A_N,B-B_N\in\mathscr C$ and
\[
 r(A-A_N)+r(B-B_N)
          \leq\sum_{j>N}(r(a_j)+r(b_j))\longrightarrow0.
 \tag{ED8}
\]
The identities with the tails hold on the common finite locus,
hence in the relative Borel quotient. Thus $h=A-B\in\mathbb D_S$
is the norm limit of the subsequence and of the original Cauchy
sequence. No abstract vector invisible to Borel capacity has been
adjoined.

For $u,v\in\mathscr C$, multilinearity gives
\[
 I(u,v)=B(u-v,u^{[q]})-B(u-v,v^{[q]}).
\]
Apply (ED4) to obtain (ED6). If a sequence in $\mathscr C$
converges in norm to $h\in\mathbb D_S$, its energy roots are
bounded by the norm triangle inequality and the preceding lemma.
Equation (ED6) makes it strongly Cauchy. R11 completeness supplies
a limit $u\in\mathscr C$. The norm limit and the strong limit
have the same local capacity class by (ED5) and the R11 strong
topology theorem. Thus $h=u$, proving norm closedness.
Convexity is the frozen cone addition property.
\end{proof}

\begin{definition}[Curvature as an energy-dual functional]
\label{def:r12-dual-curvature}
For $u\in\mathscr C$ put
\[
       J(u)(h)=B(h,u^{[q]})=-\int_Vh\,d\widehat\mu_u
                      \qquad(h\in\mathbb D_S).
 \tag{ED9}
\]
The displayed integral is literal and absolutely convergent: for
$h=a-b$, $|h|\leq(-a)+(-b)$ off a common null set, and the
two mixed first energies are finite. The current measure ignores
that exceptional set. Thus $J(u)$ is a genuine functional in
$\mathbb D_S^*$, with its original current interpretation.
\end{definition}

\begin{theorem}[Exact dual norm and strong continuity of curvature]
\label{thm:r12-dual-curvature-continuity}
For every continuous linear functional $\ell$ on $\mathbb D_S$,
\[
 \|\ell\|_*=
     \sup_{\substack{a\in\mathscr C\\r(a)\leq1}}|\ell(a)|.
 \tag{ED10}
\]
For $u\in\mathscr C$ one has $\|J(u)\|_*=r(u)^q$.
For $u,v\in\mathscr C$ and $R=\max(r(u),r(v))$,
\[
 \begin{split}
 \|J(u)-J(v)\|_*&\leq qR^{q-1}\|u-v\|_{\rm ed},\\
 \|J(u)-J(v)\|_*&\leq\sqrt q\,R^{(q-1)/2}I(u,v)^{1/2}.
 \end{split}
 \tag{ED11}
\]
At $R=0$ both sides mean zero; for $q=1$ the power $R^0$ is
one. In particular strong convergence of potentials gives norm
convergence of their source functionals, in every remaining degree.
\end{theorem}

\noindent\textbf{Status: \texttt{NEW\_PROVED}} for the precise
Banach dual and quantitative all-degree source map. The integrated
positive energy matrices used in the proof are frozen inputs.

\begin{proof}
If the supremum in (ED10) is $L$, a decomposition $h=a-b$
gives $|\ell(h)|\leq L(r(a)+r(b))$. Taking the infimum yields
$\|\ell\|_*\leq L$; the reverse follows from
$\|a\|_{\rm ed}=r(a)$. Equation (ED4) bounds the norm of
$J(u)$ by $r(u)^q$. For $r(u)>0$, evaluation on $u/r(u)$
gives equality; if $r(u)=0$, the class and functional are zero.

The finite telescoping expansion gives
\[
 (J(u)-J(v))(h)
    =\sum_{l=0}^{q-1}B(h,u-v,u^{[l]},v^{[q-1-l]}).
 \tag{ED12}
\]
Equation (ED4) proves the first bound in (ED11).
For the sharper strong-topology bound define
\[
 H(f,g)=\sum_{l=0}^{q-1}
                 B(f,g,u^{[l]},v^{[q-1-l]}).
\]
This is a positive semidefinite bilinear form on $\mathbb D_S$.
Indeed each summand is the integrated energy form with the
fixed completed positive closed mixed background. Positivity on
finite differences follows from R11 (GE7), passed through the
global mixed-energy convergence (IC13). It is not a pointwise
Alexandrov--Fenchel assertion. We have
\[
 H(u-v,u-v)=I(u,v),\qquad
 H(a,a)\leq q r(a)^2R^{q-1}\quad(a\in\mathscr C).
\]
Cauchy--Schwarz for this positive semidefinite form and (ED12)
bound the functional on every $a$ with $r(a)\leq1$.
Equation (ED10) proves the second bound. Strongly convergent
rows have bounded energy roots by R11, so the final assertion
follows. Every construction used the intrinsic original products;
no special tropical trace was substituted for their test sheaf.
\end{proof}

\paragraph{Two different source distances.}
When signed measure integration defines the corresponding continuous
functionals, their difference norm is
\[
 \delta_*(\mu,\nu)
     =\sup_{\substack{a\in\mathscr C\\r(a)\leq1}}
                    \left|\int_V(-a)\,d(\mu-\nu)\right|.
 \tag{ED13}
\]
This is not the R11 absolute-variation distance (VF11), which has
$\int(-a)\,d|\mu-\nu|$ inside the supremum. The latter dominates
(ED13) whenever it is defined, and can be strictly stronger.
The signed cancellation in (ED13) is essential to (ED11).
The Banach space and its curvature map do not prove an analytic
PS/PL presentation quotient, or equality on additional GK delta tests.
